\documentclass[../../../include/open-logic-section]{subfiles}

\begin{document}

\olfileid{fol}{syn}{its}

% Part: first-order-logic
% Chapter: semantics
% Section: intro-semantics

\olsection{引言}

赋予表达式意义属于语义学的范畴。
语义学的核心概念是在结构中的\emph{满足}关系。
结构为语言的基本构件赋予意义：一个\emph{论域}是一个非空的对象集合。
量词被解释为遍历该论域，常项符号被指派以论域中的元素，函数符号被指派以从论域到自身的函数，谓词符号被指派以论域上的关系。
论域连同对这些基本词汇的指派，共同构成一个\emph{结构}。
公式中可能出现变元，为了给出其语义，我们还必须将论域中的元素指派给它们——这便是\emph{变元指派}。
最后，\emph{满足关系}将这些部分结合起来。
一个公式可以在结构~$\Struct{M}$ 中相对于变元指派~$s$ 被满足，记作 $\Sat{M}{!A}[s]$。
这一关系同样是通过对~$!A$ 的结构进行归纳来定义的，例如，利用逻辑联结词的真值表，依据 $!A$ 和~$!B$ 是否被满足来定义 $(!A \land !B)$ 的满足条件。
由此可知，如果公式~$!A$ 是一个语句（即没有自由变元），变元指派就是无关紧要的，因此我们可以直接谈论语句在结构中是否被满足。

基于语句的满足关系 $\Sat{M}{!A}$，我们可以进而定义有效性、语义后承和可满足性这些基本的语义概念。
一个语句是\emph{有效的}，记作 $\Entails !A$，当且仅当每个结构都满足它。
一组语句 $\Gamma$ \emph{语义地后承}一个语句，记作 $\Gamma \Entails !A$，当且仅当每个满足~$\Gamma$ 中所有语句的结构也都满足~$!A$。
而一组语句是\emph{可满足的}，当且仅当存在某个结构同时满足其中的所有语句。
由于公式是归纳定义的，而满足关系又是基于公式的结构归纳定义的，因此我们可以使用归纳法来证明语义的各种性质，并建立起所定义的各个语义概念之间的联系。

\end{document}